\begin{figure}[!htbp]
\centering
\begin{tikzpicture}[node distance=8mm and 9mm,
block/.append style={text width=63mm,minimum height=16mm},
process/.append style={text width=63mm,minimum height=16mm}]
\node[block,text width=116mm] (src) at (0,0) {88 source participants: 17,419 complete four-second windows\\First screen: 17,391 approved and 28 flagged};
\node[block] (p1) at (-3.8,-2.4) {Stage 02 primary\\12,396 windows, 69 people\\4,132 windows per diagnosis};
\node[block] (left) at (3.8,-2.4) {Stage 02 leftover\\5,023 windows retained\\Includes the 28 first-screen flags};
\node[process,below=of p1] (car) {Controlled-primary CAR and QC\\65 signal exclusions; 61 balance trims\\12,270 retained (4,090 per class)};
\node[process,below=of left] (rejoin) {Later full-cohort analysis\\Uniform CAR, QC and boundary guards\\on both Stage 02 partitions};
\node[block,below=of car] (v2) {V2 boundary-safe primary\\11,942 windows, 69 people\\328 boundary-intersecting windows removed};
\node[block,below=of rejoin] (full) {V3 / V4 full cohort\\16,824 windows, 88 people\\Per-person pooling and class weights};
\draw[arrow] (src.south) -- ++(0,-0.3) -| (p1.north);
\draw[arrow] (src.south) -- ++(0,-0.3) -| (left.north);
\draw[arrow] (p1)--(car);
\draw[arrow] (car)--(v2);
\draw[arrow] (left)--(rejoin);
\draw[arrow] (p1.east) -- ++(0.4,0) |- (rejoin.west);
\draw[arrow] (rejoin)--(full);
\end{tikzpicture}
\caption{Dataset lineage. The controlled-primary path and the later full-cohort path create separate analysis indexes. Both Stage 02 partitions are preserved.}
\label{fig:dataset-flow}
\end{figure}
